\section{Desglose de preguntas y respuestas clave: de las frases de la entrevista al juicio de ingeniería}

\subsection{Sobre el límite de escala: por qué él considera que ``3 billones son posibles''}

En la entrevista, Lex pregunta directamente sobre el límite de escala, y la respuesta de Jensen no es un optimismo emocional, sino un regreso al modelo de suministro. Su lógica es la siguiente: la producción de NVIDIA no recae en una sola planta, sino que la sostiene en conjunto una cadena de suministro extensa, por lo que el límite teórico de expansión no lo determina una sola empresa. El punto central de esta respuesta no está en la cifra concreta de valor de mercado, sino en que ``el límite de la capacidad lo determina el sistema''.

Si se traslada esta misma lógica a la construcción de una plataforma de IA empresarial, la conclusión se sostiene igual. Es difícil que un solo equipo complete por sí mismo toda la cadena, desde el entrenamiento del modelo hasta su puesta en producción; al final, esto depende de que el equipo de datos, el equipo de plataforma, el equipo de negocio, el equipo de cumplimiento y el equipo de infraestructura conformen en conjunto una ``cadena de suministro organizacional''. La eficiencia de la coordinación organizacional suele determinar el límite superior más que cualquier indicador técnico aislado.

\begin{knowledgebox}{Traducir un problema de mercado de capitales a un problema de ingeniería}
La respuesta de Jensen puede traducirse en tres problemas de ingeniería:
\begin{itemize}
    \item ¿Tu sistema tiene capacidad de replicación sostenida, o solo de entrega puntual?
    \item ¿Tus dependencias clave están distribuidas, o concentradas en un único cuello de botella?
    \item ¿El ritmo de tus entregas es predecible, de modo que los clientes y socios puedan incorporarlo a su planeación?
\end{itemize}
Estas tres preguntas son igual de aplicables dentro de una organización.
\end{knowledgebox}

\subsection{Sobre si «la IA sustituirá a las personas»: la definición de trabajo que él ofrece}

La manera de expresarse de Jensen en la conversación es muy propia de la ingeniería. No presenta la IA como una entidad misteriosa, sino que la define como un amplificador: algo que permite a una persona común completar más rápido tareas que antes solo un experto podía realizar de manera confiable. Este planteamiento es consistente con la idea, mencionada antes, de ``reducir la fricción para quien empieza'', y también con el posicionamiento de producto de NVIDIA, es decir, ofrecer una infraestructura que permita que la inteligencia sea producible, desplegable y mantenible.

\textit{``AI should help you become better at your work, not just replace your work.''} Esta postura implica que el eje del gobierno organizacional debe ubicarse en el ``límite de la colaboración entre humanos y máquinas'': qué tareas deben quedar bajo el juicio final de una persona, qué pasos pueden automatizarse mediante el modelo y qué etapas requieren registros de auditoría y trazabilidad.

\begin{warningbox}{Riesgo de implementación: hablar solo de la tasa de sustitución, no de la cadena de responsabilidad}
Muchos proyectos de IA fracasan no por falta de capacidad del modelo, sino porque la asignación de responsabilidades no queda clara. Si no existe un punto de entrega definido entre humanos y máquinas, ante cualquier anomalía aparece una zona vacía en la que ``el modelo no responde y la persona tampoco''. El relato de Jensen nos recuerda que la IA en producción debe diseñar al mismo tiempo el límite de capacidad y el límite de responsabilidad.
\end{warningbox}

\subsection{Sobre el estilo de ejecución: de dónde viene la velocidad}

El estilo de ejecución que aparece en la entrevista puede resumirse como ``sincronización de alta frecuencia + concurrencia entre dominios + corrección rápida''. Este patrón parece costoso, pero cuando la complejidad del sistema es alta y los cambios externos son rápidos, en realidad reduce el costo total de comunicación. Esto se debe a que los reportes en cadena van perdiendo contexto de manera continua, mientras que la discusión concurrente permite completar la verificación de varios dominios en una sola reunión.

En las palabras de Jensen hay una señal de ejecución muy marcada: cuando un problema es suficientemente complejo, la tarea de quien dirige no es solo tomar decisiones, sino también construir un mecanismo que haga ``visible el problema''. Si el problema no es visible, ni el mejor experto puede intervenir a tiempo; si el problema es visible y puede compartirse, la organización puede entonces aprovechar su inteligencia colectiva.

\subsection{Resumen de la sección}

Este capítulo tradujo en juicios de ingeniería las frases más recurrentes de la entrevista: el límite de escala depende de la coordinación del sistema, la colaboración entre humanos y máquinas requiere un diseño de cadena de responsabilidad, y la velocidad de ejecución proviene de un mecanismo de verificación concurrente. Estos tres puntos conforman en conjunto el método operable que subyace al relato de NVIDIA.

